\subsection{Inercont}
\todo{insert whats about and how will you write this passage}
\textcolor{green}{procedure}
explain what contrastive learing is.
1 whats a Neural network: whats are the differences between deep learning and statistical models?
    a General informations
        where do we use it 
        what are these tasks
        why do we use them
        Learning Paradigms
        Types of NN
        Feet forward vs feet backward networks
        advantages disadvantages 
        Differences between statistics and deep learning (k-NN defined as Machine Learning)

    b Details: 
        Arcitecture: Layers, Nodes, activation function, bias,
        Nodes, defining, which number to choose,
        Weights, definig what happens to them
        Bias, defining, why do you need bias
        Activtion function, defining, hat type of af are there, differences
        Loss function, defining, differences, 
        Batch
        learning rate



Methods: 
Minimum Coavariance Determinant is asymptotically normal and robust. But it is hard to compute.


Therefor the authors constructed an algorithm, thats much faster than the first algorithm.
MCD Alternative: for dataset with very few observations/ excessive noise use Minimum Regularized Covariance Determinant
Insert: the curse of dimensionality

For the end: how can I be sure, that my results are legit?

for dataset analysis: multicolinearityt?



What is the solution?
The Minimum Covariance Determinant, first published in 1985 by , is a robust estimator of the covariance matrix in multidimensional dataset.
Given a dataset of n observations and a predetermined subset sie h, the MCD searches for the h observations whose empirical covariance matrix has the smalltest possible determinant.
A determinant is a mathematical property of a matrix.

How does it work?


What are the advantages?

What are the disadvantages?


\subsection{Data_Analytics:}
The requirements are met.
The Null is H0: 


   
            - Multivariate Normality: Henze-Zirkler test - if data is non normal, statistical methods will fail
            - Modality: Clustering Tendency (Hopkins Statistics) - If data has multiple naturl clusters ,multimodel, global outlier methods will fail 
        - Cardinality and Sparsity:
            - Identify the number of unique values in categorical - High-cardinality categorical variables (like "Zip Code") can create "sparse" 
                regions that look like outliers but are actually just rare valid entries.
            - Visualisation: 
        - Bivariate Analysis
            - Bridge the gap between individual data points and multicolinearity
            - Multicollinearity: VIF
            - Mutual information: identify non-linear dependencies thats standard correlation might miss
       - Visualisation: A table of features, type (numerical/categorical) and cardinality.


    Geomatric structure
        high dimensional data rearly is uniformly
        - Intrinstic dimensionality (ID): 
            - Definition: ID is the minimum number of variables (or dimensions) required to represent the essential information within datatset and without significant loss.
                is the variance just drivenby a small number of data
            - Problem: Curse of Dimensionality. The volume of the space increases so fast that the data becomes sparse
                If your dataset has 50 variables but 95\% of the variance is captured by 3 variables, the "curse of dimensionality" is low, and simple methods (like Mahalanobis Distance) might beat complex ones (like Isolation Forest).
            - Visualization: Heatmap, Use a method like the Scree Plot from PCA
        - Manifold structure: Is the data on a curved surface or is it clustered
            - Definition: manifold structure refers to the hypothesis that high dimensional data actually lives on a mich simpler, lower-dimensional surface embedded within that high-dimensional space
            - Problem: Traditional distances fail on manifold. If two points are on the opposide side of a fold, they might look close in 3D space, but they are actually very far apart along the surface of the paper.
            - Impact: A data point miht look like an outlier in 3D becaus its away from the center, but is perfectly normal when you follow the flow of the manifold. 
                or the opposite. A point could be inside thedata cloud but odd the manifold surface - making a hidden outlier
            - Analyze: Estimate Intrinsic Dimensionality using Maximum Likelihood Estimation of intrinsic dimension
            - Visualization: t-SNE plot or UMAP, those are specifically designed to unfold the manifold so you can see 2D structure 
            - Result: if the manifold are varying in density, a global method (T-Score) will fail, Manifold aware methods are more suitable (LOF) 
            
            
    Preprocessing and transformation
        - Multicollinearity: PCA
        - High cardinality: Noralization


\subsection{Experiments}
Structure: 10 - 12 pages
    Introduction
       motivate why you are using benchmarking as a indicator and then why you are applying the mothods to a real world data.
       describe your hypothesis
    Experimental Setup  
        MCD: specify size of subset 
        k-NN: specify K
        Intercont: describe the F, G networks, wich activation function, wich loss function, arcitecure
        Wich Hardware, Windows version, python version, python packages


Robustness 
    Global Anomalies 
    Find the maxmimum and minimum of the dataset and create points that are k standard deviasions away
Local Anomalies
    Take a real datapoint and movi it just far enough so its no longer in a dense cluster, but still in within the general data cloud
Cluster Anomalies
    Generate a small set of points that are very close to each other but far from the main data mass
Dependency Anomalies
    It testsif the model understands the relationship between features. Swap features values bewteen two different normal rows
    If Row A is a "High Income / High Spending" person and Row B is a "Low Income / Low Spending" person, swap their spending. 
    You now have a "High Income / Low Spending" person (which might be normal) and a "Low Income / High Spending" person (a dependency anomaly)

\subsection{Hyperparameter tuning and Sensitivity analysis}
I should do it once, because If I change the hyperparameters for every anomalie, I will test the adaptability of the method, not the robustness.
In this section I will analyse how robust and stable the results are, by changing the hyperparameters in the methods.
It's important, that here is just written the methodology. 
So I will concentrate on the how not the why. 
Therefor it is a part of the experimental design not the result section.

- Introduction 
    - Motivation: Why is it imporant to test the sensitivity?
    - Robustness: 
    - Stability: 
- MCD
    - Whats are the hyperparameter - support fraction (h)
    - contamination parameter ?
- k-NN  
    Hyperparameter: k 
    Hyperparametertuning: Maybe use Octuna as a pyhton framework
- InterCont
    Hyperparameter: u 

\subsection{k-NN: Choice of hyperparameter k, hyperparameter tuning}
\subsubsection{k-NN: Choice of hyperparameter k, hyperparameter tuning}
I tried to reproduce the F1-Score for all datasets except Abalone, KDDCUP99-Rev and KDDCUP99, because where not in the official ODDS Dataset. 
In the original Dataset there were sort of the same datasets, but I think the auther updated it. So I rejected them from further analysis. 
First I tried reproducing by using the classic Python machine leaning library sckit-learn. 
I noticed that the result differ from the ones in the paper.
As the auther of the paper noticed, they used the python library PyOD, wich is well known in the scientif community.
After implementing the code, I recognized that again the result differ substantially. 
Unfortunately in the reproduction package supplied by the authers, there are no description of reproducting k-NN.
They concentrated on their own .
For the F1 score of k-NN I suppose that they optimized for the best contmination.
Contamination is PyOD specific hyperparameter, that manages the sensitivity.
For every contamination rate = x, take x\% of the highest anomaliescore and mark it as one.
The difference to 100 are inliers.
Therefor I decided to use  ROCAUC wich is not influenced by contamination. 

\subsection{Breakdown-Point h of FAST-MCD}
\subsection{InterCont: hyperparameter k, m, u, hyperparameter tuning,}
In paper it is described how they did it. They call it ablation analysis.
First they switched the tanh activation function to LeakyReLU.
Then they just normalized the query and the vectors once, then they deployed a variant that just did the second normalization (the conventional way).
a variant in wich the order of the two normalizations is reversed. 
A variatnt with no normalizations. 


\subsection{Literatur}
\citep[p.5]{hilal2022financial}


\subsection{Metrics}
\textcolor{green}
{Comparing the anomaly detection performance of unsupervised anomaly detection algorithms
is not as straight forward as in the classical supervised classification case. In contrast to simply
compare an accuracy value or precision/recall, the order of the anomalies should be taken into
account. In classification, a wrongly classified instance is for sure a mistake. This is different in
unsupervised anomaly detection. For example, if a large dataset contains ten anomalies and
they are ranked among the top-15 outliers, this is still a good result, even if it is not perfect. To
this end, a common evaluation strategy for unsupervised anomaly detection algorithms is to
rank the results according to the anomaly score and then iteratively apply a threshold from the
first to the last rank. This results in N tuple values (true positive rate and false positive rate),
which form a single receiver operator characteristic (ROC). Then, the area under the curve
(AUC), the integral of the ROC, can be used as a detection performance measure. A nice inter-
pretation of the AUC is also given when following the proof from [70] and transform it into
the anomaly detection domain: The AUC is then the probability that an anomaly detection
algorithm will assign a randomly chosen normal instance a lower score than a randomly cho-
sen anomalous instance. Hence, we think the AUC is a perfect evaluation method and ideal for
comparison. However, the AUC only takes the ranking into account and completely neglects
the relative difference of the scores among each other. Other measures can be used to cope
with this shortcoming by using more sophisticated rank comparisons. Schubert et al. [38] com-
pares different rank correlation methodologies, for example Spearman’s pi and Kendall’s tau as an
alternative to AUC with a focus on targeting outlier ensembles. A second possible drawback of
using AUC might be that it is not ideal for unbalanced class problems and methods like area
under precision-recall curve or Matthews correlation coefficient could possibly better empha-
size small detection performance changes. Nevertheless, AUC based evaluation has been
evolved to be the de facto standard in unsupervised anomaly detection, most likely due to its
practical interpretability, and thus also serves as the measure of choice in our evaluation.
Please note that the AUC, when it is used in a traditional classification task, typically
involves a parameter, for example k, to be altered. In unsupervised anomaly detection, the
AUC is computed by varying an outlier threshold in the ordered result list. As a consequence,
if a parameter has to be evaluated (for example different k), this yields to multiple AUC values.
Another important question in this context is how to evaluate k, the critical parameter for most
of the nearest-neighbor and clustering-based algorithms. In most publications, researchers
often fix k to a predefined setting or choose “a good k” depending on a favorite outcome. We
believe that the latter is not a fair evaluation, because it somehow involves using the test data
(the labels) for training. In our evaluation, we decided to evaluate many different k’s between
10 and 50 and finally report the averaged AUC as well as the standard deviation. In particular,
for every k, the AUC is computed first by varying a threshold among the anomaly scores as
described above. Then, the mean and standard deviation for all these AUCs is reported. This
procedure basically corresponds to a random-k-picking strategy within the given interval,
which is often used in practice when k is chosen arbitrarily. The lower bound of 10 was chosen
because of statistical fluctuations occurring below. For the upper bound, we set a value of 50
such that it is still suitable for our smaller datasets. We are aware, that one could argue to
increase the upper bound for the larger datasets or even make it smaller for the small datasets
like breast-cancer. Fig 8 shows a broader evaluation of the parameter k illustrating that our
lower and upper bound is in fact useful.}

\textcolor{green}
{Although unsupervised anomaly detection does not utilize any label information in practice,
they are needed for evaluation and comparison. When new algorithms are proposed, it is com-
mon practice that an available public classification dataset is modified and the method is com-
pared with the most known algorithms such as k-NN and LOF. There is a set of typically used
datasets for classification, which are retrieved from UCI machine learning repository [61]. The
typical preprocessing comprises of selecting one class as the anomalous class and sub-sampling
some small amount of instances from that randomly. Unfortunately, the resulting datasets are
hardly published and cannot be regenerated by other scientists. Since the number of anomalies
is typically very low, a different subset might result in very different detection scores
We also put a
strong emphasis on a semantic background such the evaluation makes sense. This includes the
two assumptions that (1) anomalies are rare and (2) are different from the norm. Additionally,
we consider only point anomaly detection tasks as meaningful datasets for benchmarking since
a different preprocessing might again lead to non-comparable results.
With our dataset selection, we cover a
broad spectrum of application domains including medical applications, intrusion detection,
image and speech recognition as well as the analysis of complex systems. Additionally, the
datasets cover a broad range of properties with regard to dataset size, outlier percentage and
dimensionality. To our knowledge, this is the most comprehensive collection of unsupervised
anomaly detection datasets for algorithm benchmarking. As already stated, we published the
datasets to encourage researchers to compare their proposed algorithms with this work and
hope to establish an evaluation standard in the community \citep[p.20]{goldstein2016comparative}}
